\documentclass[11pt,a4paper]{article}

\usepackage[utf8]{inputenc}
\usepackage[T1]{fontenc}
\usepackage{amsmath,amssymb}
\usepackage{booktabs}
\usepackage{graphicx}
\usepackage{geometry}
\usepackage{hyperref}
\usepackage{multirow}
\usepackage{caption}
\usepackage{subcaption}
\usepackage{xcolor}
\usepackage{float}

\geometry{margin=2.5cm}

\title{Comparative Study of Deep Learning Architectures for\\ Atypical Lipomatous Tumor Segmentation in MRI}
\author{}
\date{\today}

\begin{document}
\maketitle

% ============================================================
\section{Introduction}
% ============================================================

Atypical lipomatous tumors (ALTs) are low-grade soft tissue neoplasms that pose diagnostic and therapeutic challenges in clinical practice. Accurate delineation of tumor boundaries from MRI is essential for treatment planning and follow-up assessment. This report presents a comparative evaluation of four deep learning-based segmentation approaches applied to ALT segmentation from T1- and T2-weighted MRI sequences.

We evaluate two architectural families---the encoder--decoder U-Net with a pretrained backbone (via \texttt{segmentation\_models\_pytorch}, SMP) and the nnU-Net-style PlainConvUNet---each in two input configurations: standard 2D (single slice) and 2.5D (three consecutive slices stacked as channels). All models are trained with 5-fold patient-level cross-validation and assessed on a held-out test set.

% ============================================================
\section{Methods}
% ============================================================

\subsection{Dataset and Preprocessing}

MRI volumes (T1 and T2) were converted to 2D axial slices. Expert-annotated binary masks delineate the tumor region. A patient-level split reserves a held-out test set; the remaining patients are partitioned into 5 folds for cross-validation. All images are resized to $256 \times 256$ pixels.

A \texttt{slice\_tumor\_map} flags slices containing tumor. During training, \emph{progressive negative sampling} linearly ramps the fraction of tumor-free (negative) slices from 0 to $\text{NEG\_SAMPLE\_RATIO} = 0.5$ across epochs, mitigating class imbalance while gradually introducing harder examples.

Data augmentation (training only) includes horizontal/vertical flips, rotations, elastic transforms, brightness/contrast jitter, and coarse dropout via the \texttt{Albumentations} library. For 2.5D models, each input is a 3-channel stack of the previous, current, and next axial slice, normalized with ImageNet statistics.

\subsection{Model Architectures}

Table~\ref{tab:architectures} summarizes the four approaches.

\begin{table}[H]
\centering
\caption{Model architecture summary.}
\label{tab:architectures}
\begin{tabular}{lllcc}
\toprule
\textbf{ID} & \textbf{Architecture} & \textbf{Encoder / Config} & \textbf{Input Ch.} & \textbf{Params} \\
\midrule
M1 & SMP U-Net 2D     & EfficientNet-B4 (ImageNet) + scSE decoder & 1 & 20.3\,M \\
M2 & SMP U-Net 2.5D   & EfficientNet-B4 (ImageNet) + scSE decoder & 3 & 20.3\,M \\
M3 & nnU-Net 2D       & PlainConvUNet, 7 stages, deep supervision  & 1 & 29.5\,M \\
M4 & nnU-Net 2.5D     & PlainConvUNet, 7 stages, deep supervision  & 3 & $\sim$29.5\,M \\
\bottomrule
\end{tabular}
\end{table}

\paragraph{SMP U-Net (M1, M2).}
Built with \texttt{segmentation\_models\_pytorch}. The encoder is an EfficientNet-B4 pretrained on ImageNet. The decoder uses spatial and channel squeeze-and-excitation (scSE) attention. Transfer learning is employed: the encoder is frozen for the first 5 epochs, then unfrozen with a 10$\times$ reduced learning rate for fine-tuning.

\paragraph{nnU-Net PlainConvUNet (M3, M4).}
Uses the \texttt{dynamic\_network\_architectures} implementation of the nnU-Net plain convolutional U-Net. The network has 7 stages with feature maps $[32, 64, 128, 256, 320, 512, 720]$, 2 convolutions per stage, and stride-2 downsampling. Deep supervision is enabled with geometrically decaying weights across resolution levels. Unlike M1/M2, no pretrained weights are used---the model is trained from scratch.

\subsection{Training Protocol}

All models share the following training configuration:

\begin{itemize}
    \item \textbf{Optimizer:} AdamW, initial learning rate $3 \times 10^{-4}$
    \item \textbf{Scheduler:} Warmup cosine annealing (3 warmup epochs, 50 total epochs)
    \item \textbf{Batch size:} 8
    \item \textbf{Loss:} Combined loss = $0.8 \cdot \mathcal{L}_\text{Dice} + 0.15 \cdot \mathcal{L}_\text{Focal} + 0.05 \cdot \mathcal{L}_\text{Boundary}$
    \item \textbf{Early stopping:} Patience of 10 epochs on validation Dice
    \item \textbf{Gradient clipping:} Max norm 1.0
    \item \textbf{Mixed precision:} Enabled on CUDA
    \item \textbf{Cross-validation:} 5-fold, patient-level splits
\end{itemize}

\subsection{Inference Enhancements}

At test time, two optional enhancements are evaluated:
\begin{itemize}
    \item \textbf{Test-Time Augmentation (TTA):} The input is passed through the model in four orientations (original, horizontal flip, vertical flip, both flips). The resulting logits are averaged before applying sigmoid activation.
    \item \textbf{Post-Processing (PP):} Morphological operations and connected-component filtering to remove small spurious predictions.
\end{itemize}

\subsection{Evaluation Metric}

The primary metric is the \textbf{Dice Similarity Coefficient (DSC)}, defined as:
\[
\text{DSC} = \frac{2 |P \cap G|}{|P| + |G|}
\]
where $P$ is the predicted segmentation and $G$ is the ground truth.

% ============================================================
\section{Results}
% ============================================================

\subsection{Best Validation Dice per Fold}

Table~\ref{tab:val_dice} reports the best validation Dice achieved in each fold during training.

\begin{table}[H]
\centering
\caption{Best validation Dice per fold and cross-validation summary.}
\label{tab:val_dice}
\begin{tabular}{lccccccc}
\toprule
\textbf{Model} & \textbf{Fold 1} & \textbf{Fold 2} & \textbf{Fold 3} & \textbf{Fold 4} & \textbf{Fold 5} & \textbf{Mean} & \textbf{Std} \\
\midrule
M1: SMP U-Net 2D    & 0.6625 & 0.6851 & 0.7211 & \textbf{0.8552} & 0.7473 & 0.7342 & 0.0671 \\
M2: SMP U-Net 2.5D  & 0.6246 & 0.6166 & 0.7354 & 0.8305 & 0.7663 & 0.7147 & 0.0828 \\
M3: nnU-Net 2D      & 0.6631 & 0.6935 & 0.6249 & 0.8301 & 0.5802 & 0.6784 & 0.0848 \\
M4: nnU-Net 2.5D    & 0.4546 & 0.6772 & 0.6606 & 0.7476 & 0.5525 & 0.6185 & 0.1031 \\
\bottomrule
\end{tabular}
\end{table}

\begin{table}[H]
\centering
\caption{Best validation loss per fold.}
\label{tab:val_loss}
\begin{tabular}{lccccccc}
\toprule
\textbf{Model} & \textbf{Fold 1} & \textbf{Fold 2} & \textbf{Fold 3} & \textbf{Fold 4} & \textbf{Fold 5} & \textbf{Mean} & \textbf{Std} \\
\midrule
M1: SMP U-Net 2D    & 0.5469 & 0.5317 & 0.3725 & \textbf{0.4305} & 0.3981 & 0.4559 & 0.0707 \\
M2: SMP U-Net 2.5D  & 0.5632 & 0.5400 & 0.3343 & 0.4706 & 0.3928 & 0.4602 & 0.0866 \\
M3: nnU-Net 2D      & 0.5921 & 0.5447 & 0.4041 & 0.4735 & 0.4591 & 0.4947 & 0.0662 \\
M4: nnU-Net 2.5D    & 0.6415 & 0.5824 & 0.3912 & 0.5308 & 0.4621 & 0.5216 & 0.0880 \\
\bottomrule
\end{tabular}
\end{table}

\paragraph{Observations.}
Across all models, Fold~4 consistently achieves the highest validation Dice (0.75--0.86), suggesting that particular patient partition is more favorable. M1 (SMP U-Net 2D) achieves the highest mean validation Dice at $0.7342 \pm 0.0671$, followed by M2 (SMP U-Net 2.5D) at $0.7147 \pm 0.0828$. M3 (nnU-Net 2D) follows at $0.6784 \pm 0.0848$, while M4 (nnU-Net 2.5D) ranks last at $0.6185 \pm 0.1031$---the lowest mean and highest variance of all models, indicating that combining 2.5D input with training from scratch is particularly challenging.

\subsection{Test Set Evaluation}

Table~\ref{tab:test_dice} reports test Dice for each fold's best model, both without and with TTA + post-processing.

\begin{table}[H]
\centering
\caption{Test Dice per fold --- without and with TTA + post-processing (TTA+PP).}
\label{tab:test_dice}
\begin{tabular}{lcccccc}
\toprule
\textbf{Model} & \textbf{Fold 1} & \textbf{Fold 2} & \textbf{Fold 3} & \textbf{Fold 4} & \textbf{Fold 5} & \textbf{Mean} \\
\midrule
\multicolumn{7}{l}{\textit{Without TTA+PP}} \\
\midrule
M1: SMP U-Net 2D    & \textbf{0.8060} & 0.7241 & 0.7781 & 0.6979 & 0.6683 & 0.7349 \\
M2: SMP U-Net 2.5D  & 0.7274 & 0.7578 & 0.6997 & 0.7123 & 0.6323 & 0.7059 \\
M3: nnU-Net 2D      & 0.7783 & 0.7650 & 0.5983 & 0.7669 & 0.5803 & 0.6978 \\
M4: nnU-Net 2.5D    & 0.6480 & 0.7501 & 0.6599 & 0.6656 & 0.7032 & 0.6854 \\
\midrule
\multicolumn{7}{l}{\textit{With TTA+PP}} \\
\midrule
M1: SMP U-Net 2D    & 0.8212 & 0.7800 & 0.8056 & \textbf{0.8362} & 0.7284 & \textbf{0.7943} \\
M2: SMP U-Net 2.5D  & 0.7782 & 0.7836 & 0.8049 & 0.7523 & 0.7033 & 0.7645 \\
M3: nnU-Net 2D      & 0.8066 & 0.8238 & 0.6810 & 0.8233 & 0.6305 & 0.7530 \\
M4: nnU-Net 2.5D    & 0.7793 & 0.7914 & 0.7413 & 0.7923 & 0.6559 & 0.7520 \\
\bottomrule
\end{tabular}
\end{table}

\subsection{Reproducibility: SMP U-Net 2.5D Second Run}

The SMP U-Net 2.5D model (M2) was trained a second time to assess reproducibility. Table~\ref{tab:m2_runs} compares both runs.

\begin{table}[H]
\centering
\caption{SMP U-Net 2.5D (M2): comparison of two independent training runs.}
\label{tab:m2_runs}
\begin{tabular}{lcccc}
\toprule
& \multicolumn{2}{c}{\textbf{Validation Dice}} & \multicolumn{2}{c}{\textbf{Test Dice (TTA+PP)}} \\
\cmidrule(lr){2-3} \cmidrule(lr){4-5}
\textbf{Fold} & \textbf{Run 1} & \textbf{Run 2} & \textbf{Run 1} & \textbf{Run 2} \\
\midrule
Fold 1 & 0.6246 & 0.6404 & 0.7782 & 0.7484 \\
Fold 2 & 0.6166 & 0.7052 & 0.7836 & 0.6998 \\
Fold 3 & 0.7354 & 0.7520 & 0.8049 & 0.7189 \\
Fold 4 & 0.8305 & 0.8569 & 0.7523 & 0.5909 \\
Fold 5 & 0.7663 & 0.7875 & 0.7033 & 0.5145 \\
\midrule
Mean $\pm$ Std & $0.7147 \pm 0.083$ & $0.7484 \pm 0.073$ & $0.7645$ & $0.6545$ \\
\bottomrule
\end{tabular}
\end{table}

Run~2 achieves marginally higher validation Dice ($0.7484$ vs.\ $0.7147$) but substantially worse test Dice ($0.6545$ vs.\ $0.7645$). This gap suggests potential overfitting in Run~2, where higher validation performance does not generalize to unseen test data. The discrepancy highlights the importance of evaluating on a truly held-out set rather than relying solely on cross-validation metrics.

\subsection{Summary Comparison}

Table~\ref{tab:summary} presents the aggregated performance metrics.

\begin{table}[H]
\centering
\caption{Overall comparison of all approaches (completed models).}
\label{tab:summary}
\begin{tabular}{lccc}
\toprule
\textbf{Model} & \textbf{Val Dice (mean $\pm$ std)} & \textbf{Test Dice} & \textbf{Test Dice (TTA+PP)} \\
\midrule
M1: SMP U-Net 2D     & $0.7342 \pm 0.0671$ & 0.7349 & \textbf{0.7943} \\
M2: SMP U-Net 2.5D   & $0.7147 \pm 0.0828$ & 0.7059 & 0.7645 \\
M3: nnU-Net 2D       & $0.6784 \pm 0.0848$ & 0.6978 & 0.7530 \\
M4: nnU-Net 2.5D     & $0.6185 \pm 0.1031$ & 0.6854 & 0.7520 \\
\bottomrule
\end{tabular}
\end{table}

% ============================================================
\section{Discussion}
% ============================================================

\subsection{SMP U-Net vs.\ nnU-Net}

The SMP U-Net models (M1, M2) consistently outperform the nnU-Net models (M3) in both validation and test Dice. This advantage is attributable to:
\begin{enumerate}
    \item \textbf{Pretrained encoder:} EfficientNet-B4 weights from ImageNet provide strong low- and mid-level features, enabling faster convergence and better generalization on small medical datasets.
    \item \textbf{scSE attention:} The spatial and channel squeeze-and-excitation blocks in the decoder help recalibrate feature maps, improving boundary precision.
    \item \textbf{Encoder freezing strategy:} Freezing the encoder initially prevents catastrophic forgetting of pretrained features, while gradual unfreezing allows domain-specific fine-tuning.
\end{enumerate}

The nnU-Net models, despite having 45\% more parameters (29.5M vs.\ 20.3M), must learn all features from scratch on a relatively small dataset. Deep supervision helps gradient flow but does not fully compensate for the lack of pretrained features. M4 (nnU-Net 2.5D) in particular exhibits the worst validation Dice ($0.6185 \pm 0.1031$) with the highest variance, though its test performance with TTA+PP (0.7520) nearly matches M3 (0.7530), suggesting that post-processing can partially compensate for weaker learned representations.

\subsection{2D vs.\ 2.5D Input}

Counter to expectations, the 2.5D variants do not improve over their 2D counterparts:
\begin{itemize}
    \item M1 (2D) outperforms M2 (2.5D) by $+0.030$ in test Dice with TTA+PP ($0.7943$ vs.\ $0.7645$).
    \item M3 (2D) and M4 (2.5D) achieve nearly identical test Dice with TTA+PP ($0.7530$ vs.\ $0.7520$), but M3 has substantially better validation Dice ($0.6784$ vs.\ $0.6185$), suggesting more stable training.
\end{itemize}

Possible explanations include:
\begin{itemize}
    \item The inter-slice spacing may be too large for adjacent slices to provide useful contextual information.
    \item The 2.5D approach triples the effective input dimensionality without adding proportional spatial coherence, potentially diluting informative features with noise from adjacent slices that may not contain tumor.
    \item With a small dataset, the 2.5D approach may introduce additional variance that hinders generalization.
\end{itemize}

\subsection{Impact of TTA and Post-Processing}

TTA + post-processing consistently improves test Dice across all models:
\begin{itemize}
    \item M1: $+0.059$ ($0.7349 \to 0.7943$)
    \item M2: $+0.059$ ($0.7059 \to 0.7645$)
    \item M3: $+0.055$ ($0.6978 \to 0.7530$)
    \item M4: $+0.067$ ($0.6854 \to 0.7520$)
\end{itemize}

The improvement ranges from 5.5 to 6.7 percentage points across all models, confirming that TTA+PP is architecture-agnostic and should be adopted as standard practice. Notably, M4 benefits the most from TTA+PP, likely because its weaker base predictions have more room for improvement through augmentation-based ensembling.

\subsection{Fold Variance and Generalization}

A notable observation is the high variance across folds. Fold~4 consistently achieves the best validation Dice (0.83--0.86) across all models, while Fold~5 tends to perform worst. This likely reflects the heterogeneity of the patient population and the small dataset size, where particular patient combinations can strongly influence fold-level metrics.

The M2 reproducibility experiment further underscores this: despite random seed control via \texttt{KFold(random\_state=42)}, stochastic training dynamics produce meaningfully different test outcomes between runs.

% ============================================================
\section{Conclusion}
% ============================================================

Among the four evaluated approaches, \textbf{M1 (SMP U-Net 2D)} with test-time augmentation and post-processing achieves the best overall performance with a mean test Dice of \textbf{0.7943}. Key findings:

\begin{enumerate}
    \item \textbf{Pretrained encoders matter:} SMP U-Net with ImageNet-pretrained EfficientNet-B4 outperforms the nnU-Net trained from scratch, despite the latter having more parameters and deep supervision.
    \item \textbf{2.5D input does not help:} Stacking adjacent slices as channels does not improve segmentation quality in this setting, possibly due to inter-slice spacing or dataset size limitations.
    \item \textbf{TTA+PP is essential:} Test-time augmentation with post-processing provides a consistent $\sim$6\% Dice improvement at no training cost.
    \item \textbf{Small dataset challenges:} High fold-level variance and the reproducibility gap in M2 highlight the need for larger datasets or more robust regularization strategies.
\end{enumerate}

Future work should investigate (i) 3D volumetric architectures for true spatial context, (ii) self-supervised or foundation model pretraining on medical data, (iii) ensemble strategies combining the best fold models, and (iv) larger or multi-institutional datasets to reduce fold-level variance.

\end{document}
